% ch-codeplan --- the code and test plan (operational companion to ch13)

\Cref{ch:roadmap} says \emph{what} gets built and in which phase. This chapter says \emph{how} each piece is
tested before it is trusted. It is the engineering contract behind the roadmap's gates.
\Cref{chcp:sec:harness} fixes the four test layers and what a green result in each layer actually certifies.
\Cref{chcp:sec:forward} states the test plan as frozen pass/fail protocols, each mapped to the roadmap phase
that owns it. \Cref{chcp:sec:uq} summarizes the uncertainty-quantification methods and the production
checklist. \Cref{chcp:sec:ci} closes with the continuous-integration gates and the definition of done, tied
to the decision registry of \cref{ch13:sec:registry}. No results appear in this chapter: every protocol below
is planned, not run. Synthetic-data verdicts will be recorded in \cref{ch:results-synth}; real-data verdicts
will be recorded in \cref{ch:results-spain}. Evaluation-metric definitions are not restated below: the
evaluation constitution of \cref{ch07:sec:constitution} is the sole authority on what ``calibrated'' and
``better'' mean.

\section{Testing philosophy and the four-layer harness}
\label{chcp:sec:harness}

A calibration code base fails in two distinct ways. Either the code computes the wrong number, or the code
computes the right number for a procedure that cannot recover the truth. Conventional software tests catch
only the first failure. We therefore run four layers of tests. Each layer certifies something the previous
layer cannot, and the roadmap's phase gates (\cref{ch13:sec:reading}) are defined in terms of them.

\paragraph{Layer 1: unit tests on hand-computable micro-examples.} Every nontrivial data transformation has a
companion example small enough to verify by hand. Three anchor cases: (i) a three-firm choice set with known
funding histories, where the pools, the excluded most-recently-funded firm, the newcomer row, and the winner
labels produced by \code{build\_choice\_sets} are asserted against a worked-out table; (ii) a two-bucket
compensator for an exponential-kernel intensity, checked against the closed-form integral to machine
precision; (iii) deduplication on a constructed deal table containing exact duplicates, near-duplicates a few
days apart, and legitimate consecutive-week repeats, asserting which rows survive. These tests certify
arithmetic and indexing, nothing more. But when a recovery test fails, they let us clear the plumbing in
minutes.

\paragraph{Layer 2: equivalence tests between fast and reference fitters.} The repository ships slow,
transparent loop fitters and vectorized fast fitters (\code{fit\_sector\_glm\_fast}, \code{fit\_choice\_fast}).
The fast versions exploit structure: the sector GLM is separable per receiving sector, so twelve small concave
problems replace one large one; the ranker is a padded-tensor batch over events and candidates. Speed is
worthless if the optimum moves. On subsamples, both implementations must agree at the optimum --- same loss to
$10^{-6}$, parameters to matching tolerance --- before a fast fitter may be used in any experiment, and the
check is rerun whenever either implementation changes. Every new likelihood must also pass
analytic-vs-numeric gradient checks at the same fixed tolerances.

\paragraph{Layer 3: statistical recovery tests.} These are the scientific tests. We fit on synthetic data
whose ground truth we wrote down, and a run passes when the truth lies within 2 standard errors of the
estimate across 5 seeds. A recovery failure with passing layers 1--2 is a finding about the \emph{method} ---
identifiability, confounding, misspecification. Findings of this kind will be recorded in
\cref{ch:results-synth} and feed the decision registry. Recovery tests are the only layer that can certify the
phrase ``the pipeline works,'' which is why every roadmap gate requires a recovery rehearsal before a new
fitter touches real data.

\paragraph{Layer 4: invariant and leakage audits.} Two standing assertions run over the whole feature stack.
The point-in-time assertion recomputes every feature using only data truncated at time $t$ and requires the
result to be bit-identical to the production value: any divergence means some path read the future. The
oracle-isolation assertion requires that no fitting path opens \code{ground\_truth.npz} or
\code{private\_deal\_map.csv}; oracle files are for evaluation only, and a fit that touches them is void
regardless of its metrics. Both audits are structural: they do not check that a number is right, they check
that it was legal to compute. On real data this layer becomes \modrepo{data/audits.py} of Phase 0
(\cref{ch13:sec:phase0}), with the oracle role played by vendor-side fields that postdate the prediction time.

\begin{table}[htbp]
\centering
\small
\begin{tabular}{lp{5.6cm}l}
\toprule
Layer & What it certifies & Cadence \\
\midrule
Unit (micro-examples) & arithmetic and indexing of each transformation & every commit \\
Equivalence (fast vs.\ loop) & the vectorized optimum is the reference optimum & every commit \\
Recovery (synthetic truth) & the procedure can find truth within 2 SE, 5 seeds & nightly, small configs \\
Invariant / leakage & features are point-in-time; oracle files untouched & blocking, every commit \\
\bottomrule
\end{tabular}
\caption{The four test layers and what a green result actually means.}
\label{chcp:tab:layers}
\end{table}

Three reproducibility rules bind everything above. First, every reported number must be regenerable from one
seeded entry point --- for the rehearsal worlds themselves,
\begin{verbatim}
PYTHONPATH=. python -m synthetic.generate --regime A --seed 0
\end{verbatim}
and analogously one script per experiment table with a \code{--seed} argument. Second, temporal splits only,
everywhere: train weeks $[0,416)$ against test weeks $[416,520)$ on the weekly track, an 80\% day quantile on
the choice track, and no random split anywhere in the stack, including model selection --- this is the same
temporal discipline the evaluation constitution imposes (\cref{ch07:sec:temporal}). Third, runtimes are logged
alongside results, because a fitter that silently becomes ten times slower is a regression even when its
numbers do not move.

\begin{codehooks}
Layer 1--2 suites live under \code{tests/}. Recovery campaigns are the seeded entry points under
\code{experiments/}, whose verdicts will be transcribed into \cref{ch:results-synth}; the rehearsal stage
itself is \code{synthetic/} (\cref{ch13:sec:registry}). On real data, layer 4 becomes \modrepo{data/audits.py}
(P0.2) and the consolidated metric harness is \modrepo{eval/} (P2.4); both replace per-experiment ad hoc
scripts.
\end{codehooks}

\section{The forward test plan}
\label{chcp:sec:forward}

What follows is the committed test plan. Each protocol states its purpose, its procedure in a few lines, a
pass criterion frozen \emph{before} the run so that results can fail, and the roadmap phase that owns its
real-data form. Every protocol has status \emph{planned}; this chapter reports no outcomes. Synthetic runs
write their verdicts into \cref{ch:results-synth}; real-data reruns write theirs into \cref{ch:results-spain}.
Unless stated otherwise, every experiment runs on both rehearsal regimes at 5 seeds with error bars, under the
splits and audit rules of \cref{chcp:sec:harness}, scored only by metrics the constitution admits
(\cref{ch07:sec:constitution}). \Cref{chcp:tab:forward} summarizes.

\begin{protocol}[E1: bucket width, decay identifiability, and latent confounding]
\label{chcp:prot:e1}
\emph{Purpose:} quantify what fine-grained timestamps buy over coarse buckets, and how badly latent
confounding contaminates fitted excitation. \emph{Procedure:} fit the weekly GLM and a multivariate event-time
sector Hawkes (exponential kernel, log-linear macro baseline) on the same streams; rerun the event-time fit at
1d, 7d, and 28d bucket widths, reporting bias and SE of the branching scale $\branch$ and decay $\decay$ at
each width; separately, regenerate regime A with latent-factor strength in $\{0, 0.15, 0.3\}$ at fixed true
radius and plot fitted effective radius against latent strength, with and without mitigations (lagged
aggregate count covariate, off-diagonal penalties). \emph{Pass:} $\decay$ recovered within 2 SE at day
resolution on regime B while the $\branch$--$\decay$ ridge visibly reopens by 28d; the confounding curve is
monotone and yields a stated bucket-width recommendation. \emph{Status:} planned; results to
\cref{ch:results-synth}. \emph{Phase:} the real-data descendant is the honest-resolution decision of Phase 0
(\cref{ch13:sec:phase0}, D6.1) and the ridge diagnostics standing on the Phase 2 Hawkes fit
(\cref{ch13:sec:phase2}).
\end{protocol}

\begin{protocol}[E2: choice-model sub-experiments]
\label{chcp:prot:e2}
\emph{Purpose:} settle the four open design choices inside the conditional logit before the ranker is frozen
for production. \emph{Procedure and pass criteria, per sub-experiment:}
\begin{enumerate}[label=E2.\arabic*,leftmargin=3.2em]
\item \emph{Gap functional form.} Fit log-gap versus binned gap dummies ($<$1m, 1--3m, 3--6m, 6--12m, $>$12m)
versus exponential decay $\exp(-\text{days}/\tau)$ with $\tau \in \{30, 90, 180\}$ days; select by held-out
NLL and plot each implied hazard-versus-gap curve against truth. Pass: the binned fit recovers the true
26-week step in regime A and the smooth exponential decay in regime B.
\item \emph{Newcomer designs N1--N4.} Compare the bare ASC against the context-covariate ASC (sector heat and
macro features), the representative-entrant construction (plug-in and integrated Gaussian variants), the
nested two-step, and the immigrant/offspring split that aligns first financings with the Hawkes background
stream. Pass: the recommended design calibrates quarterly newcomer share within 2 percentage points, with
incumbent weights materially unchanged. The design space is owned by the outside-option construction of
\cref{ch:universe}.
\item \emph{Sampled-softmax stability.} Refit at maximum candidate counts $\{32, 64, 128\}$ and with the full
denominator on a 10\% subsample. Pass: weights stable across candidate counts, consistent with the
subsampled-denominator consistency result of \cref{app:proofs}.
\item \emph{Exclusion-rule sensitivity.} Vary the exclude-most-recent rule (sector-scoped, global, none);
count dropped repeats and measure weight shifts. Pass: the effect is pinned with a number, and the rule is
stated in the production spec either way.
\end{enumerate}
\emph{Status:} planned, all four; results to \cref{ch:results-synth}. \emph{Phase:} E2.1, E2.3, and E2.4
belong to the Phase 1 linear ranker (\cref{ch13:sec:phase1}); E2.2 rides the Phase 2 outside-option work
(\cref{ch13:sec:phase2}).
\end{protocol}

\begin{protocol}[E3: two-stage factorization versus single-layer block Hawkes]
\label{chcp:prot:e3}
\emph{Purpose:} the central architecture question --- given exact dates and truthful risk sets, what does the
two-stage factorization still buy, and what does it cost in joint-intensity coherence? \emph{Procedure:} fit
the full firm-level event-time block Hawkes (log-link baseline, self-inhibition, size-normalized sector field)
on day-resolution timestamps; evaluate it and the two-stage pipeline on identical held-out events for joint
NLL, next-deal sector prediction, and funded-firm ranking conditional on the event. \emph{Pass:} on regime B
the single-layer fit recovers the self-inhibition ordering with rank correlation at least 0.9; the
head-to-head produces a written recommendation with numbers on all three metrics, whichever way it falls.
\emph{Status:} planned; results to \cref{ch:results-synth}. \emph{Phase:} behind the Phase 3 trigger for
signed continuous-time kernels (\cref{ch13:sec:phase3}) --- it runs only if a consuming question demands a
structural continuous-time claim that routes (a)/(c) cannot express, and its recovery rehearsal comes first.
\end{protocol}

\begin{protocol}[E4: neural-operator placement]
\label{chcp:prot:e4}
\emph{Purpose:} establish whether and where the neural-operator surrogate for the compensator equations pays
for itself, given that exact likelihoods are the committed default (\cref{ch:censoring,ch:gof}).
\emph{Procedure:} report forward-solver accuracy versus sector count $M$ and the wall-clock crossover against
the exact solvers under repeated forward solves; benchmark differentiable-solver inverse calibration against
the MLE for bias, variance, and wall-clock on identical data. \emph{Pass:} the forward solver reaches a
sub-2\% relative $L^2$ band; the crossover point (in solves and in $M$) is documented; the inverse-calibration
table states plainly whether the MLE wins at this scale. \emph{Status:} planned; results to
\cref{ch:results-synth}. \emph{Phase:} not on the critical path; revisited under Phase 3 rules only if a
repeated-forward-solve workload (scenario simulation in \cref{ch:decision}) materializes.
\end{protocol}

\begin{protocol}[E5: walk-forward backtest with covariate ablation]
\label{chcp:prot:e5}
\emph{Purpose:} the end-to-end rehearsal of the production loop under strictly forward information.
\emph{Procedure:} freeze the best when/where and which-firm models; walk forward weekly over the final 20\% of
days, emitting sector intensities, ranked candidate lists, and newcomer probabilities; score top-1/5/10 and
MRR against oracle and random, PIT count calibration with the quasi-Poisson dispersion correction, and lead
time (weeks a funded firm spends in the top decile before its deal); rerun with macro covariates removed and
on the covariate-scale 1.5 world. \emph{Pass:} the pipeline beats random throughout and degrades gracefully
toward oracle; the covariate ablation turns ``covariates help'' into a signed number in both directions.
\emph{Status:} planned; results to \cref{ch:results-synth}. \emph{Phase:} its real-data form is the Phase 2
gate itself --- the end-to-end weekly run producing the deliverables of \cref{ch:decision}
(\cref{ch13:sec:phase2}), driven by the rolling-origin machinery of \modrepo{eval/}.
\end{protocol}

\begin{protocol}[B1: boosted-ranker benchmark]
\label{chcp:prot:b1}
\emph{Purpose:} a model-free discriminative ceiling for the identity layer. \emph{Procedure:} train XGBoost
rankers \citep{chenguestrin2016} on the identical choice sets and splits in two variants --- feature parity
(exactly the six logit features, isolating functional form) and kitchen sink (all point-in-time-legal raw
fields, snapshot columns explicitly forbidden); score with per-event renormalization; modest hyperparameter
search with a time-ordered validation tail. \emph{Pass (frozen decision rule):} a material win for the boosted
variant means the logit is missing features or nonlinearity, and the pre-committed response is the
grouped-softmax boosted ranker of \cref{ch:ranking}, which keeps the conditional-likelihood family while
granting the score function interactions; a narrow result keeps the interpretable logit. \emph{Status:}
planned; results to \cref{ch:results-synth}. \emph{Phase:} runs against the Phase 1 linear baseline
(\cref{ch13:sec:phase1}); the grouped-softmax response, if triggered, ships in Phase 2
(\cref{ch13:sec:phase2}). Both variants are rerun on real data, with verdicts in \cref{ch:results-spain},
before any response is declared necessary there.
\end{protocol}

\begin{protocol}[B2: non-homogeneous Poisson null and improvement decomposition]
\label{chcp:prot:b2}
\emph{Purpose:} show, not assume, that self-excitation earns its complexity at this data scale.
\emph{Procedure:} fit the nested ladder --- covariates-only inhomogeneous Poisson (simulable by thinning
\citep{lewisshedler1979}), covariates plus excitation, full event-time timing --- plus a seasonal-naive
baseline; decompose held-out improvement across the rungs with residual diagnostics per rung \citep{ogata1988},
on the NLL ladder of \cref{ch07:sec:ladder}. \emph{Pass:} the decomposition exists with error bars over 5
seeds; if the excitation rung contributes approximately nothing, that is written into the registry, not
smoothed over. \emph{Status:} planned; results to \cref{ch:results-synth}. \emph{Phase:} the real-data rerun
\emph{is} the Phase 1 gate --- the nats-ladder note on the 2011--2024 extract (\cref{ch13:sec:phase1}),
recorded in \cref{ch:results-spain}.
\end{protocol}

\begin{protocol}[B3: nonparametric intensity challenger]
\label{chcp:prot:b3}
\emph{Purpose:} test whether the nonparametric intensity models of \cref{ch:nonparam} --- boosted trees,
forests, and small networks on the discretized likelihood --- beat the parametric families they challenge.
\emph{Procedure:} fit each \cref{ch:nonparam} model in the same harness as every other fitter: same temporal
splits, same audits, same 5 seeds, same feature-legality rules; compare against the log-linear GLM of
\cref{ch:poisson} and the parametric Hawkes of \cref{ch:hawkes} on held-out likelihood (the NLL ladder of
\cref{ch07:sec:ladder}) and on calibration (the PIT and dispersion checks of \cref{ch07:sec:counts}).
\emph{Pass (frozen decision rule):} the comparison table exists with error bars; a nonparametric model is
adopted only if it wins on held-out likelihood \emph{and} calibrates at least as well as the parametric
incumbent; a likelihood win with worse calibration is a finding, not an adoption. \emph{Status:} planned;
results to \cref{ch:results-synth}, real-data rerun to \cref{ch:results-spain}. \emph{Phase:} rides the
Phase 1 baseline ladder (\cref{ch13:sec:phase1}) and is rerun against the Phase 2 architecture
(\cref{ch13:sec:phase2}).
\end{protocol}

\begin{protocol}[Causal-pipeline validation on known ground truth]
\label{chcp:prot:causal}
\emph{Purpose:} test the causal stack of \cref{ch:eucausal} against a world where the estimand is known.
\emph{Procedure:} generate a treatment with known selection on observables and on the latent quality mark;
estimate naive, backdoor-adjusted, and AIPW effects; attach cluster-bootstrap confidence intervals and
E-values \citep{hernan2020}. \emph{Pass, asserted in code so a failure is a red test:} the naive estimate is
biased in the selection direction; adjustment removes the observable component; AIPW matches backdoor with
better robustness to nuisance misspecification; and the residual gap to truth is attributable to latent
quality with magnitude consistent with the computed E-value. \emph{Status:} planned; results to
\cref{ch:results-synth}. \emph{Phase:} a blocking rehearsal gate before any estimator of \cref{ch:eucausal}
runs on real data, and a precondition for the Phase 3 frailty escalations (\cref{ch13:sec:phase3}).
\end{protocol}

\begin{table}[htbp]
\centering
\footnotesize
\begin{tabular}{lp{5.6cm}p{2.2cm}l}
\toprule
Test & Pass criterion (frozen) & Status & Phase \\
\midrule
E1 bucketing & $\decay$ within 2 SE at day resolution on B; ridge visible by 28d & planned & P0/P2 \\
E1 confounding curve & monotone fitted-radius curve; mitigations quantified & planned & P2 \\
E2.1 gap form & binned fit recovers the true step (A), smooth decay (B) & planned & P1 \\
E2.2 newcomer N1--N4 & newcomer share within 2pp; incumbent weights stable & planned & P2 \\
E2.3 sampled softmax & weights stable across 32/64/128/full-subsample & planned & P1 \\
E2.4 exclusion rule & sensitivity pinned with a number & planned & P1 \\
E3 head-to-head & rank-corr $\ge 0.9$ on B; three-metric comparison written & planned & P3 (trigger) \\
E4 operator & sub-2\% forward band; crossover and inverse tables & planned & P3 (trigger) \\
E5 backtest & beats random throughout; ablation deltas signed & planned & P2 gate \\
B1 boosted ranker & frozen decision rule applied & planned & P1$\to$P2 \\
B2 Poisson null & improvement decomposition with error bars & planned & P1 gate \\
B3 nonparametric challenger & likelihood win counts only with equal-or-better calibration & planned & P1$\to$P2 \\
Causal validation & all four coded assertions pass against truth & planned & pre-\cref{ch:eucausal} \\
\bottomrule
\end{tabular}
\caption{Forward test plan at a glance. All rows: both regimes where applicable, 5 seeds, temporal splits,
audits of \cref{chcp:sec:harness} in force. Synthetic verdicts go to \cref{ch:results-synth}; real-data
verdicts go to \cref{ch:results-spain}; phases per \cref{ch:roadmap}.}
\label{chcp:tab:forward}
\end{table}

\section{Uncertainty quantification and the production checklist}
\label{chcp:sec:uq}

Every interval this document reports comes from one of five techniques, and each carries a different
guarantee. \Cref{chcp:tab:uq} is deliberately blunt about which.

\begin{table}[htbp]
\centering
\footnotesize
\begin{tabular}{p{2.9cm}p{2.9cm}p{3.4cm}p{4.2cm}}
\toprule
Method & Technique & Guarantee & Key assumptions \\
\midrule
Count GLM (\cref{ch:poisson}) & quasi-Poisson sandwich & asymptotic SE validity for mean parameters & correct mean function; dispersion as nuisance \\
Two-stage / choice (\cref{ch:ranking}) & cluster bootstrap by week & asymptotic CI coverage & weak dependence across clusters; clustering at the right level \\
Boosted prediction (\cref{ch:buckets}) & split conformal / CQR \citep{romano2019} & finite-sample marginal coverage & exchangeability; weakened under temporal splits \\
Hawkes MLE (\cref{ch:hawkes}) & Laplace / posterior interval & local asymptotic normality & correct specification; stationarity \\
MBPP (\cref{ch:censoring}) & parametric bootstrap & coverage under the fitted DGP & model correctness; report $\branch$, not $\decay$ \\
\bottomrule
\end{tabular}
\caption{UQ summary: what each interval means and what it silently assumes.}
\label{chcp:tab:uq}
\end{table}

Only the conformal row carries a finite-sample statement, and even that one is marginal (not conditional) and
degrades when temporal splits break exchangeability. Every other interval is asymptotic and
specification-dependent. One consequence is worth spelling out: if the held-out Pearson dispersion of a count
fit is $\phi$, naive Poisson standard errors are understated by a factor of about $\sqrt{\phi}$ until
corrected --- which is why the dispersion report is a checklist item, not a courtesy. The production checklist
below is the run-time counterpart of the Phase 0 audit suite (\cref{ch13:sec:phase0}): the audits verify the
data once, the checklist binds every fit and report thereafter.

\begin{enumerate}[label=C\arabic*.,leftmargin=3.0em]
\item \emph{Point-in-time audit.} Reconstruct every feature from data knowable at the prediction date; join
macro series on publish date, never reference date (\cref{ch:data}).
\item \emph{Deduplication.} Apply the validated dedup window before any counting; duplicates and date jitter
alias into spurious short-lag excitation (\cref{ch:data,ch:censoring}).
\item \emph{Risk-set truthfulness audit.} Validate the membership mask $\riskset_t$ against a labelled
subsample of known entrants and exits; a contaminated mask can invert fitted weights (\cref{ch:ranking}).
\item \emph{Effective-radius guard.} Refuse to simulate any fitted model at effective radius $\ge 0.95$; for
log-link raw-count feedback refuse well below that, or use bounded transforms (\cref{ch:hawkes}).
\item \emph{Dispersion correction.} Report held-out Pearson dispersion and scale standard errors accordingly;
dispersion far above 1 is a specification alarm, not a nuisance (\cref{ch07:sec:counts}).
\item \emph{Universe coverage and newcomer calibration.} Track the untracked-firm share and recalibrate the
newcomer ASC whenever the pool definition changes; the ASC is meaningless detached from its pool
(\cref{ch:universe,app:proofs}).
\item \emph{Temporal splits only.} No random splits anywhere --- not for model selection, not for conformal
calibration, not for the benchmarks of \cref{chcp:sec:forward} (\cref{ch07:sec:temporal}).
\item \emph{Seeds and regeneration.} Every reported number carries a seed and a regeneration command; a table
without an entry point under \code{experiments/} does not exist.
\end{enumerate}

\section{Continuous integration and the definition of done}
\label{chcp:sec:ci}

The four layers of \cref{chcp:sec:harness} map onto three CI gates. On every commit, the unit and equivalence
suites run and must pass; they are fast because the micro-examples are tiny and the equivalence checks run on
subsamples. Nightly, the recovery suite runs on reduced configurations --- fewer weeks, fewer seeds --- as an
early-warning system; a nightly recovery regression blocks the next release, not the next commit. The leakage
audit is a blocking check at both cadences: a commit that makes any feature non-point-in-time, or that lets a
fitting path touch \code{ground\_truth.npz} or \code{private\_deal\_map.csv}, does not merge, whatever its
metrics say. The registry review of \cref{ch13:sec:registry} sits on top of all three: a pull request that
contradicts a decision box is rejected or arrives with a documented reversal, and CI cannot green its way
around that rule.

The artifact policy is equally plain. Results tables are versioned CSVs, written by the experiment entry
points and diffed like code, so a silently shifted number is a visible change in review. Figures are
regenerated from those CSVs by scripts and are never hand-edited; a figure that cannot be regenerated is
deleted rather than trusted. Campaign verdicts are transcribed append-only into \cref{ch:results-synth} for
synthetic runs and into \cref{ch:results-spain} for real-data runs, so the evidentiary trail behind every
decision box stays auditable.

\begin{decision}{D-CP.1 --- the harness is constitutional}
The four-layer harness of \cref{chcp:sec:harness}, the three CI gates above, and the frozen-pass-criterion
discipline of \cref{chcp:sec:forward} are binding on all work in \modrepo{}. No fast fitter runs in an
experiment without a current layer-2 equivalence pass; no fitter touches real data without a layer-3 recovery
rehearsal on its own family; no pass criterion is edited after its experiment has run. Rationale: the value of
a test plan is precisely that results are allowed to fail, and these rules make failure visible and durable.
Reversal trigger: none anticipated; amendments go through the registry (\cref{ch13:sec:registry}).
\end{decision}

\begin{decision}{D-CP.2 --- definition of done}
The test campaign is done when three conditions hold. (i) Every row of \cref{chcp:tab:forward} has moved from
\emph{planned} to \emph{pass} or \emph{fail} with a number attached --- not a narrative, a number --- and the
verdict recorded in \cref{ch:results-synth}. (ii) The benchmarks \cref{chcp:prot:b1,chcp:prot:b2,chcp:prot:b3}
are run on the real extract, recorded in \cref{ch:results-spain}, and cited by the phase gates of
\cref{ch13:sec:phase1,ch13:sec:phase2}, so the architecture choices rest on measured evidence from the data
they will serve, not on rehearsal priors alone. (iii) Every failed criterion is either fixed and rerun, or
documented as a limitation in the registry --- because a plan whose failures disappear quietly was never a
test plan at all. Rationale: this is the operational meaning of the roadmap's gates. Reversal trigger: a
change in \cref{ch:decision}'s consuming actions that adds or retires a protocol, per D7.1.
\end{decision}
